\documentclass{article}
\usepackage[a4paper, total={6.5in, 10in}]{geometry}
\setlength{\parindent}{0pt}
\usepackage{tcolorbox}
\tcbset{colback=white!94!black, colframe=black!60!white, coltitle = white, boxrule=0.8pt, arc=2mm}
\newtcolorbox{boxs}[2][]{title= #2,#1}
\usepackage{datetime2}
\usepackage[british]{babel}
\usepackage{amsfonts}
\usepackage{amsmath}

\title{\textbf{Linear Algebra HW 7}}
\author{Robert O'Connell}
\date{\today}
\begin{document}
\maketitle
\vspace{5mm}
\subsection*{Question 1}
\subsubsection*{(a)}
\[
\chi_A(\lambda) = \det(A - \lambda I)
\]
\[
\chi_A(\lambda) = 0 \iff (-5 - \lambda)(-5- -\lambda) = 0
\]
Thus we see that \(\lambda = -5\) is the only eigenvalue of \(A\)
\\

\(\mathcal{N}(A- 5 I)\):
\[
\begin{bmatrix}
    0 & 0 & | & 0 \\
    5 & 0 & | & 0 
\end{bmatrix}
\]
The dimension of this is \(1\), thus there exists only one linearly independent 
eigenvector of \(A\), e.g. \(\begin{bmatrix}
    0 \\
    1
\end{bmatrix}\)

Thus we know that the Jordan form must be \(\begin{bmatrix}
    5 & 1 \\
    0 & 5
\end{bmatrix}\)

\subsubsection*{(b)}
Now looking in the next nullspace
\(\mathcal{N}(A - 5I)^2\):
\[
\begin{bmatrix}
    0 & 0 & | & 0 \\
    0 & 0 & | & 0
\end{bmatrix}
\]
One can clearly see that this space is all of \(\mathbb{R}^2\), picking one of these 
such that is linearly independent from our earlier eigenvector, we get 
\(\begin{bmatrix}
    1 \\ 
    0
\end{bmatrix}\)

This set \(\mathcal{C} = \left\{\begin{bmatrix}
0 \\ 1
\end{bmatrix}, \begin{bmatrix}
1 \\ 0
\end{bmatrix} \right\} \) then forms a Jordan basis for \(A\)

Thus \(M =  P_{\mathcal{C} \rightarrow \mathcal{E}}\begin{bmatrix}
    0 & 1 \\
    1 & 0
\end{bmatrix}\)


\subsection*{Question 2}
\subsubsection*{(a)}
\[
\chi_A(\lambda) = \det(A - \lambda I)
\]
\[
\chi_A(\lambda) = 0 \iff (4-\lambda)((-1-\lambda)(-4 - \lambda) - 3) -3(7(-4-\lambda) +24) + 4(7 + 8(-1-\lambda))
\]
\[
\iff - \lambda^3 - \lambda^2 + 8 \lambda +12 = 0
\]
\[
\iff (\lambda -3)(\lambda +2)^2 = 0
\]
Thus we see that \(\lambda = -2\) and \(\lambda = 3\) are the two eigenvalues of \(A\)
\\

\(\mathcal{N}(A - 3I)\):
\[
\begin{bmatrix}
    1 & 3 & 4 & | & 0  \\
    7 & -4 & 3 & | & 0 \\
    -8 & 1 & -7 & | & 0 
\end{bmatrix}
\]
\[
 \Rightarrow \begin{bmatrix}
    1 & 0 & 1 & | & 0 \\
    0 & 1 & 1 & | & 0 \\
    0 & 0 & 0 & | & 0 
\end{bmatrix}
\]
The set of eigenvectors associated with \(\lambda = 3\) is then \(\left\{\begin{bmatrix}
    -t \\ -t \\ t
\end{bmatrix}\mid t \in \mathbb{R} \right\}\)
\\

\(\mathcal{N}(A + 2I)\):
\[
\begin{bmatrix}
    6 & 3 & 4 & | & 0 \\
    7 & 1 & 3 & | & 0 \\
    -8 & 1 & -2 & | & 0 
\end{bmatrix}
\]
\[
\Rightarrow \begin{bmatrix}
    1 & 0 & \frac{1}{3} & | & 0 \\
    0 & 1 & \frac{2}{3} & | & 0 \\
    0 & 0 & 0 & | & 0
\end{bmatrix}
\]
The set of eigenvectors associated with \(\lambda = -2\) is then \(\left\{\begin{bmatrix}
    -\frac{t}{3} \\ -\frac{2t}{3} \\ t
\end{bmatrix} \mid t \in \mathbb{R} \right\}\)
\\

\(\mathcal{N}(A + 2I)^2\)
\[
\begin{bmatrix}
    25 & 25 & 25 & | & 0 \\
    25 & 25 & 25 & | & 0 \\
    -25 & -25 & -25 & | & 0 
\end{bmatrix}
\]
\[
\Rightarrow \begin{bmatrix}
    1 & 1 & 1 & | & 0 \\
    0 & 0 & 0 & | & 0 \\
    0 & 0 & 0 & | & 0  
\end{bmatrix}
\]
The set of  generalised eigenvectors associated with \(\lambda = -2\) is then \(\left\{\begin{bmatrix}
    -s -t \\ s \\ t
\end{bmatrix} | s,t \in \mathbb{R} \right\}\)
\\

Thus the generalised eigenspace associated with \(\lambda = 3 \) is \( V(3) = \left\{\begin{bmatrix}
    -t \\ -t \\ t
\end{bmatrix}\mid t \in \mathbb{R} \right\}\)
\\

and the generalised eigenspace associated with \(\lambda = -2\) is 
\( V(-2) = \left\{\begin{bmatrix}
    -s -t \\ s \\ t
\end{bmatrix} | s,t \in \mathbb{R} \right\}
\)
\\

Note: We do not need to include \(\mathcal{N}(A + 2I)\) in the second eigenspace since it is contained within
the set of generalised eigenvectors

\subsubsection*{(b)}
Since \(A\) has two eigenvectors associated with \(\lambda = 3\) and \(\lambda = -2\), and one 
generalised eigenvector associated with \(\lambda = -2\), we know that the Jordan form of \(A\)
must be (up to permuations of the blocks)
\[
J = \begin{bmatrix}
    3 & 0 & 0 \\
    0 & -2 & 1 \\
    0 & 0 & -2
\end{bmatrix}
\]

\subsubsection*{(c)} 
Picking one eigenvector associated with \(\lambda = 3\) such as \(\begin{bmatrix}
    -1 \\ -1 \\ 1
\end{bmatrix}\), 
then one generalised eigenvector associated with \(\lambda = - 2\)
such that it is linearly independent of the set of eigenvectors associated with \(\lambda = -2\) 
such as \(\begin{bmatrix}
    -2 \\ 1 \\ 1
\end{bmatrix}\), then our final vector in the Jordan basis is \((A + 2I)\) multiplied to the generalised eigenvector
which comes out to be \(\begin{bmatrix}
    -5 \\ -10 \\ 15
\end{bmatrix}\)
\\

These three vectors then form a Jordan basis

Then we know that \(M\) has as its columns the images of these basis vector, hence
\[
M = \begin{bmatrix}
    -1 & -5 & -2 \\
    -1 & -10 & 1 \\
    1 & 15 & 1
\end{bmatrix}
\]


\subsection*{Question 3}
To determine the Jordan form of \(A\), we must first look in the null spaces of 
\((A - 5I)^k, \ k \in \mathbb{N}\)
\\

\(\mathcal{N}(A - 5I)\):
\[
\begin{bmatrix}
    0 & -3 & -2 & | & 0 \\
    1 & -5 & -3 & | & 0 \\
    -1 & 8 & 5 & | & 0 
\end{bmatrix}
\]
\[
\begin{bmatrix}
    1 & -5 & -3 & | & 0 \\
    0 & -3 & -2 & | & 0 \\
    0 & 0 & 0 & | & 0 
\end{bmatrix}
\]

We see that \(\dim(\mathcal{N}(A - 5I)) = 1\), since there is only one free variable,
thus there is only linearly independent eigenvector of \(A\), thus the Jordan form of 
\(A\) must take the form 
\[
J = \begin{bmatrix}
    5 & 1 & 0 \\
    0 & 5 & 1 \\
    0 & 0 & 5
\end{bmatrix}
\]
since any other Jordan form would require at least two eigenvectors


\subsection*{Question 4}
We know every Jordan matrix is of the form \(J = M^{-1} A M\), where \(M\) is an 
invertible matrix

Thus \(A  = MJM^{-1}\)
\\

Picking \(M = \begin{bmatrix}
    1 & 0 & 0 \\
    0 & 1 & 1 \\
    0 & 0 & 1
\end{bmatrix}\), we get
\[
A = \begin{bmatrix}
    1 & 0 & 0 \\
    0 & 1 & 1 \\
    0 & 0 & 1
\end{bmatrix}
\begin{bmatrix}
    -7 & 1 & 0 \\
    0 & -7 & 0 \\
    0 & 0 & 1
\end{bmatrix}
\begin{bmatrix}
    1 & 0 & 0 \\
    0 & 1 & -1 \\
    0 & 0 & 1
\end{bmatrix}
\]
\[
A = \begin{bmatrix}
    -7 & 1 & -1 \\
    0 & -7 & 8 \\
    0 & 0 & 1
\end{bmatrix}
\]


\subsection*{Question 5}
Since the matrix \(A\) is a \(3 \times 3\) matrix with three distinct eigenvalues 
we know that it's Jordan form will be \(\begin{bmatrix}
    2 & 0 & 0 \\
    0 & 5 & 0 \\
    0 & 0 & 0 
\end{bmatrix}\)
\\

Also \(A\) is of the form \(A = MJM^{-1}\) where the columns of \(M\) are the eigenvectors
of \(A\) with respect to the standard basis vectors, 

i.e. \(M = \begin{bmatrix}
    2 & 0 & 3 \\
    3 & -1 & -1 \\
    1 & 0 & 2
\end{bmatrix}\)
\\

Then \(M^{-1} = \begin{bmatrix}
    2 & 0 & -3 \\
    7 & -1 & -11 \\
    -1 & 0 & 2
\end{bmatrix}\)

Finally,
\[
A = \begin{bmatrix}
    2 & 0 & 3 \\
    3 & -1 & -1 \\
    1 & 0 & 2
\end{bmatrix} \begin{bmatrix}
    2 & 0 & 0 \\
    0 & 5 & 0 \\
    0 & 0 & 0 
\end{bmatrix}
\begin{bmatrix}
    2 & 0 & -3 \\
    7 & -1 & -11 \\
    -1 & 0 & 2
\end{bmatrix}
\]
\[
A = \begin{bmatrix}
    8 & 0 & -12 \\
    -23 & 5 & 37 \\
    4 & 0 & 6
\end{bmatrix}
\]


\subsection*{Question 6}
Let \(T\) and \(S\) be linear operators defined on the vector space \(V\),
such that \(S \circ T = T \circ S\)
\\

Let \(\vec{k} \in \ker(S)\ \quad \) i.e. \(S(\vec{k}) = \vec{0}\)

Then since \(T\) is a linear operator we have that \(T(S(\vec{k})) = \vec{0}\)
and hence \(S(T(\vec{k})) = \vec{0}\)
\\

i.e. \(T(\vec{k}) \in \ker(S)  \ \forall \ k \in \ker(S)\)

Hence \(\ker(S)\) is invariant under \(T\)


\subsection*{Question 7}
Let \(\mathcal{B} = \{1, x, x^2, x^3, x^4\}\) be a basis for \(\mathbb{P}_4[x]\)

Then \([T]_{\mathcal{B}, \mathcal{B}}\) is the matrix which has as its columns the images of these basis vectors
\[
[T]_{\mathcal{B},\mathcal{B}} = \begin{bmatrix}
    0 & 0 & 0 & 0 & 0 \\
    0 & 1 & 0 & 0 & 0 \\
    0 & 0 & 2 & 0 & 0 \\
    0 & 0 & 0 & 3 & 0 \\
    0 & 0 & 0 & 0 & 4
\end{bmatrix}
\]

Since \([T]_{\mathcal{B},\mathcal{B}}\) is a diagonal matrix, we know that
the eigenvalues of \(T\) are the entries along the diagonal, \(0,1,2,3,4\)
and the eigenvectors are the spans of each of the vectors in \(\mathcal{B}\)


\end{document}